% Generated corrected-evidence block; see registry and audit.
\subsection{Day-in-the-Life Operations}\label{sec:sim-leop}
\paragraph{Simulation context and assumptions.} The corrected LEOP 24-hour, No-LEOP 24-hour and LEOP 365-day studies form the operational evidence used here. Their contact and SAA timestamps identify an epoch of 1 March 2000 UTC; the approved 1 January 2029 mission configuration was not simulated by these runs. Each scenario is assessed separately. Producing configuration records do not specify RAAN, mass, battery capacity, integrator, hardware selection or elevation mask. The day cases share recorded positions and contact windows, but their comparison is a scenario sensitivity rather than an isolated causal test of LEOP.

Telemetry occupancy and energy integrals hold each state over $[t_i,t_{i+1})$ and exclude terminal-row duration. Transition and slew-event logs use differing duration definitions, detailed in Appendix~\ref{app:ditl-detail}. Battery percentages use the 25~Wh planning reference. Acquisition gates and expected counts use the study's legacy response; they describe observing opportunity, not accepted events from the current FPDU. These assumptions apply to all DITL results in this section.

\begin{table}[H]\centering\small
\caption{Integrated DITL results using left-held telemetry intervals}\label{tab:annual-ditl}
\begin{tabularx}{\linewidth}{@{}Xrrr@{}}\toprule
Metric & LEOP 24 h & No-LEOP 24 h & LEOP 365 d\\\midrule
Recorded timestamp span (s) & 86,400 & 86,400 & 31,536,000\\
Cold and Dark (s) & 122 & 0 & 122\\
Detumble (s) & 2,460 & 0 & 2,460\\
Sun Pointing (s) & 79,268 & 81,852 & 29,660,918\\
Detection (s) & 2,676 & 2,674 & 894,054\\
Relay (s) & 894 & 894 & 269,876\\
Slewing (s) & 980 & 980 & 375,688\\
Safe (s) & 0 & 0 & 332,882\\
Geographic SAA (s) & 2,696 & 2,696 & 881,882\\
Acquisition gate (s) & 758 & 758 & 247,258\\
Minimum recorded battery (Wh) & 15.128 & 18.004 & 15.098\\
Recorded solar integral (Wh) & 145.575 & 150.403 & 50,288.809\\
Recorded load integral (Wh) & 85.167 & 85.719 & 30,866.103\\
Native generation (source MB) & 0.009 & 0.009 & 3.514\\
Native transfer (source MB) & 0.009 & 0.009 & 3.514\\
\bottomrule\end{tabularx}\end{table}
\begin{figure}[H]\centering\includegraphics[width=\linewidth]{evidence/science-ditl/figures/corrected_20261109_day_occupancy_battery.png}
\caption{Day-run mode occupancy and battery energy normalized to the 25 Wh planning reference.}\label{fig:annual-mode-exposure}\end{figure}
\paragraph{Startup and science acquisition.} Both day runs record 2,696~s of geographic SAA dwell and 758~s of acquisition gate. The LEOP case completes one 2,460~s detumble event, followed by science and Relay operations within the first day. Removing startup does not increase the sampled acquisition total in this pair. The simulation therefore supports the evaluated startup sequence without identifying startup as the limiting factor for that day's acquisition opportunity.

\paragraph{Annual operational continuity.} The annual timeline allocates 29,660,918~s to Sun Pointing, 894,054~s to Detection, 269,876~s to Relay and 375,688~s to Slewing, while recording 332,882~s in Safe Mode. It demonstrates a full-year schedule containing science, communications and recovery states. Safe Mode occupancy is a real operational interruption; it is retained in the resource and exposure accounting. The acquisition gate totals 247,258~s against 881,882~s geographic SAA dwell, or 28.038\%. This establishes usable modeled acquisition opportunity, with geographic dwell and Detection occupancy distinguished from actual gated exposure. Independent passes and calibrated science-tier success require the detector-response assessment in Section~\ref{payload-subsystem}.

\paragraph{Power behavior.} Both day cases and the annual case retain positive recorded battery energy, with minima of 15.128, 18.004 and 15.098~Wh respectively. Annual recorded generation of 50,288.809~Wh exceeds the modeled load integral of 30,866.103~Wh. Together with the battery trajectory, these results support sustained operation of the evaluated load schedule. The energy-integral difference is not a stored-energy or component margin; charging limits, losses and selected-cell limits must be represented before interpreting it as such.

\subsection{Ground Station and Data Return}\label{sec:sim-coverage}
The corrected annual contact log contains 820 intervals totaling 244,016 s. Mean daily contact is 668.537 s and minimum daily contact 514 s; the longest intercontact gap is 48,458 s. Visibility alone corresponds to 585,638,400 bytes at 19.2 kbps. Intersecting these contacts with telemetry-labelled Relay yields 178,682 s, or 428,836,800 ideal bytes. The transition-log calculation also gives 428,836,800 bytes of eligible capacity, despite differing mode-duration definitions. The recurring opportunities support scheduled ground relay for this evaluated orbit.

Native annual transfer credit is 3.514083862 source-labelled MB, separate from hypothetical v3 delivery. The source-labelled MB units are retained without assuming a byte conversion. The v3 preferred EMSC workload generates 46,637,364 bytes with peak modeled queue 75,336 bytes and final backlog 0 bytes. It uses corrected, time-aligned sampled gates and Relay visibility at the full ideal 2400 bytes/s; it is not emitted v3 telemetry. Appendix~\ref{app:science-data-budget} gives tier and outage sensitivity.
The calculated service accommodates the sampled annual workload and returns the queue to zero. This supports the scheduled data-return concept under the full ideal gross-rate assumption. RF margin and actual net throughput remain separate from this capacity result. The seven-day outage case requires 1,010,028 bytes of buffering, establishing the need for persistent spill beyond the candidate 256~KiB RAM queue.
